\section{Self-Correction: Dirichlet Prediction Falsified}
\label{sec:auto_correction_dirichlet}
\label{sec:auto_correction}

The K4 duality (\S\ref{sec:dualite_higgs_zeta}) suggests a natural
extension to primitive Dirichlet $L$-functions via Haar
factorisation $\bmod\,q$. This prediction was tested and falsified.
This section honestly documents the test and draws from it the
reformulation of the status of the residue $1/8$.

\subsection{The prediction P3}
\label{sec:P3}

If K4 is structurally extensible, the generic formula for the
Maslov phase of $L(s, \chi \bmod q)$ would be
\begin{equation}
\Delta_N(L \bmod q) \;=\; \sper^{2} \cdot \mathrm{Haar}(\Integers/q\Integers) \cdot \dfrac{1}{N_{\mathrm{branches}}} \;=\; \dfrac{1}{8q}.
\label{eq:P3}
\end{equation}
Predictions: $q = 3 : 1/24$, $q = 5 : 1/40$, $q = 7 : 1/56$,
$q = 11 : 1/88$.

\subsection{Test M1 (mpmath, $T \le 200$, 4 points per $q$)}
\label{sec:M1}

First test, inconclusive: signal/noise too low (4 points,
$\sigma_{\mathrm{SE}} \sim 0.2$, target signal $\sim 0.04$).

\subsection{Test M1bis (LMFDB, $T \le 200$, ${\sim} 140$ zeros per $q$)}
\label{sec:M1bis}

In-depth test on pre-computed \texttt{LMFDB} zeros for primitive
quadratic characters $q \in \{3, 5, 7, 11\}$ ($a$ parity of the
character):

\begin{table}[ht]
\centering
\small
\begin{tabular}{@{}rrrrrrrr@{}}
\toprule
$q$ & $a$ & $N$ zeros & $\langle \Delta_N \rangle$ & SE & $1/(8q)$ & $|t|_{\mathrm{PT}}$ & $|t|_{0}$\\
\midrule
3 & 1 & 114 & $-2.4 \cdot 10^{-4}$ & $0.019$ & $+0.042$ & \textbf{2.19} & $0.01$\\
5 & 0 & 129 & $-2.6 \cdot 10^{-3}$ & $0.021$ & $+0.025$ & $1.34$ & $0.13$\\
7 & 1 & 140 & $+1.8 \cdot 10^{-3}$ & $0.021$ & $+0.018$ & $0.77$ & $0.09$\\
11 & 1 & 155 & $-2.8 \cdot 10^{-4}$ & $0.021$ & $+0.011$ & $0.55$ & $0.01$\\
\bottomrule
\end{tabular}
\caption{Results of test M1bis. The \emph{t-stat} compares the
observed mean to the PT prediction $+1/(8q)$ (column $|t|_{\mathrm{PT}}$)
and to the null hypothesis $\Delta_N = 0$ (column $|t|_{0}$).}
\label{tab:M1bis}
\end{table}

Discriminatory $\chi^{2}$ test on 4 d.o.f.:

\begin{table}[ht]
\centering
\begin{tabular}{@{}lrr@{}}
\toprule
\textbf{Hypothesis} & $\chi^{2}$ & $p$-value\\
\midrule
$H_0$: $\Delta_N = 0$ & $0.023$ & $\approx 1.00$\\
$H_{\mathrm{PT}}$: $\Delta_N = +1/(8q)$ & $7.49$ & $0.11$\\
$H_{-\mathrm{PT}}$: $\Delta_N = -1/(8q)$ & $7.05$ & $0.13$\\
\bottomrule
\end{tabular}
\caption{$\chi^{2}$ test: $H_0$ is preferred by $\sim 300\times$
over $H_{\mathrm{PT}}$. The prediction P3 is \textbf{effectively
falsified} at $T \le 200$.}
\label{tab:chi2_dirichlet}
\end{table}

\paragraph{Sign discussion.}
On the four values of $q$ tested, PT predicts a strictly positive
shift $+1/(8q)$; we observe three negative shifts ($q = 3, 5, 11$)
and one positive ($q = 7$). The sign agreement alone is therefore
$1/4$, well below the $1$ that the prediction would impose.
Falsification is not only amplitude-wise: the sign structure itself
contradicts the prediction $+1/(8q) > 0$.

\subsection{Calibration on $\zeta$}
\label{sec:calibration_zeta}

With the same methodology on the first 50 zeros of $\zeta$
(\cite{Odlyzko1992}):
\begin{equation}
\langle \Delta_N \rangle_\zeta \;=\; +1.004 \pm 0.028.
\label{eq:calib_zeta}
\end{equation}
Comparisons:

\begin{table}[ht]
\centering
\begin{tabular}{@{}lrr@{}}
\toprule
\textbf{Target} & \textbf{Value} & $t$\\
\midrule
Pole of $\zeta$ alone ($+1$) & $1.000$ & $+0.15$ (compatible)\\
Pole + Maslov ($1 + 1/8 = 1.125$) & $1.125$ & $-4.29$ (rejected)\\
Maslov alone ($1/8$) & $0.125$ & $+31.29$ (rejected)\\
\bottomrule
\end{tabular}
\caption{Calibration on $\zeta$: the $+1$ from the pole entirely
absorbs the $1/8$ in the counting statistic.}
\label{tab:calibration_zeta}
\end{table}

\subsection{Reformulation of the status of the residue $1/8$}
\label{sec:reformulation_residu}

\paragraph{Critical consequence.}
Proposition~\ref{prop:densite_NPT} states that
$\NPT(\gamma) - N_{\mathrm{RvM}}(\gamma) = 1/8$, verified to
$10^{-15}$ over four decades. This difference is between \textbf{two
smooth asymptotic forms} ($\NPT$ and $N_{\mathrm{RvM}}$), not a
measurable residue in the exact counting of zeros. At the count of
real zeros, the $+1/8$ between PT and RvM is absorbed by the other
terms in the formula and does not appear as a measurable mean shift.
The $+1$ from the $\zeta$ pole entirely absorbs the $1/8$ in the
counting statistic.

The asymptotic identity~\eqref{eq:residu_un_huitieme} remains valid
as an \textbf{asymptotic identity between forms}, but the
translation ``$1/8$ visible in the zero count'' is erroneous. This
explains a posteriori why the prediction $1/(8q)$ has been
falsified: it attributed to Dirichlet a measurable residue analogous
to that of $\zeta$, while even for $\zeta$ this residue is not
measurable in the zero count at heights $T \sim 200$.

\subsection{Position in the programme}
\label{sec:P3_position}

The falsification of P3 does not affect the HP-PT structure. Remain
intact:
\begin{itemize}[nosep,leftmargin=*]
\item the PT spectral model of $\zeta$ zeros via $\Treg(s)$:
$17/17$ zeros captured at $|\Delta| < 1.5$ on $t \in [10, 70]$,
$74/79$ on $t \in [10, 200]$ with the raw formula, $75/79$ at
$|\Delta| < 0.5$ after explicit regularisation $R_1$ with
$\eps = 0.2$, and $267/269$ at $|\Delta| < 1.5$ extended to
$t \in [10, 500]$;
\item the cohomological identification
$1/8 = \sper^{2}/2 = c_1/N_{\mathrm{corners}} = \sper/4$
for $\zeta$;
\item the Higgs $\leftrightarrow \zeta$ duality via the $p = 2$
channel (conjecture~\ref{conj:K4}) as a structural strong conjecture
for $\zeta$.
\end{itemize}

What falls: the natural extension to Dirichlet via the
$\bmod\,q$ Haar measure. The PT mechanism described here is
\textbf{specific to $\zeta$}, not to the Dirichlet family via this
factorisation. The PT \emph{ledger} marks P3 as
\textbf{[FALSIFIED, $T \le 200$]}.

This is the programme's falsifiability discipline: a testable PT
prediction was produced by K4, tested, rejected by data, and the
programme restructures cleanly.
